\documentclass[11pt,a4paper]{article}
% preamble.tex — shared setup for human-verification documents
% Compile each document with:  tectonic <File>.tex

\usepackage[T1]{fontenc}
\usepackage[utf8]{inputenc}
\usepackage{lmodern}
\usepackage{amssymb,amsmath,amsthm}
\usepackage{mathtools}
\usepackage[dvipsnames]{xcolor}
\usepackage{listings}
\usepackage{hyperref}
\usepackage[margin=2.5cm]{geometry}
\usepackage{parskip}

% ── Colours ──────────────────────────────────────────────────────────
\definecolor{leanBg}{HTML}{F7F7F0}
\definecolor{haskellBg}{HTML}{F0F4F7}
\definecolor{keyword}{HTML}{7B3294}
\definecolor{comment}{HTML}{5A8C32}
\definecolor{noteFg}{HTML}{B35900}

% ── Lean 4 listings language ─────────────────────────────────────────
% NOTE: Lean listing blocks use ASCII approximations of unicode symbols
% (e.g. -> for arrow, o. for compose).  The literate key substitutes
% these with proper math-mode renderings in the PDF.
\lstdefinelanguage{Lean4}{
  morekeywords={def,theorem,structure,where,namespace,end,import,fun,
                let,have,show,by,exact,intro,cases,rfl,subst,obtain,
                funext,simp,rw,apply,variable,Type,Prop,instance,
                section,open,class,extends,deriving,sorry,noncomputable},
  sensitive=true,
  morecomment=[l]{--},
  morecomment=[n]{/-}{-/},
  morestring=[b]",
}

% ── Lean listing style ───────────────────────────────────────────────
\lstdefinestyle{lean}{
  language=Lean4,
  backgroundcolor=\color{leanBg},
  basicstyle=\small\ttfamily,
  keywordstyle=\color{keyword}\bfseries,
  commentstyle=\color{comment}\itshape,
  breaklines=true,
  frame=single,
  framerule=0.4pt,
  xleftmargin=4pt,
  xrightmargin=4pt,
  aboveskip=6pt,
  belowskip=4pt,
  literate=
    {->}{{$\to$}}2
    {=>}{{$\Rightarrow$}}2
    {<|}{{$\langle$}}1
    {|>}{{$\rangle$}}1
    {==>}{{$\Longrightarrow$}}3
    { o. }{{$\;\circ\;$}}3
    {forall}{{$\forall$}}1
    {:=}{{$\,:=\,$}}2
    {!=}{{$\neq$}}2
    {<=}{{$\leq$}}2
    {>=}{{$\geq$}}2,
}

% ── Haskell listing style ───────────────────────────────────────────
\lstdefinestyle{haskell}{
  language=Haskell,
  backgroundcolor=\color{haskellBg},
  basicstyle=\small\ttfamily,
  keywordstyle=\color{keyword}\bfseries,
  commentstyle=\color{comment}\itshape,
  breaklines=true,
  frame=single,
  framerule=0.4pt,
  xleftmargin=4pt,
  xrightmargin=4pt,
  aboveskip=6pt,
  belowskip=4pt,
  literate=
    {->}{{$\to$}}2
    {=>}{{$\Rightarrow$}}2,
}

% ── Theorem environments ────────────────────────────────────────────
\theoremstyle{definition}
\newtheorem{defn}{Definition}[section]
\newtheorem{thm}[defn]{Theorem}
\newtheorem{prop}[defn]{Proposition}

\theoremstyle{remark}
\newtheorem*{fidelitynote}{Fidelity note}

% ── Convenience macros ───────────────────────────────────────────────
\newcommand{\lean}[1]{\texttt{#1}}      % inline Lean identifier
\newcommand{\hask}[1]{\texttt{#1}}      % inline Haskell identifier
\newcommand{\Ext}{\mathrm{Ext}}
\newcommand{\Pos}{\mathrm{Pos}}
\newcommand{\Shape}{\mathrm{Shape}}


\title{Human Verification: \texttt{Containers.DContFunctor}\\[4pt]
  \large Covariant Functors of Directed Containers}
\author{Machine-generated from \texttt{lean/DContFunctor.lean} for review by Neil Ghani}
\date{July 2026}

\begin{document}
\maketitle

\begin{abstract}
This document presents every definition and theorem \emph{statement} from the
Lean~4 module \lean{Containers.DContFunctor} in three parallel forms: the
original Lean, a readable Haskell rendering, and precise mathematical English.
Proofs are omitted --- they are machine-checked by Lean.  The purpose is to
let a human verifier confirm that the formalisation captures the intended
mathematics.
\end{abstract}

\tableofcontents
\bigskip

%===================================================================
\section{Overview}
%===================================================================

This file defines \textbf{covariant functors of directed containers} --- the
morphism notion needed for Construction~\emph{con:sourcing} in the supply-chains
paper.  Contrast with \lean{DContMorphism} (in \lean{Containers.DContCat}),
whose position map is \emph{contravariant} and corresponds to a \emph{cofunctor}
of the induced categories.

A \lean{DContFunctor C D} is a \emph{functor} $\mathcal{C}(C) \to \mathcal{C}(D)$:
\begin{itemize}
  \item the shape map $f : S_C \to S_D$ acts on objects \emph{covariantly};
  \item the position map $f^\flat_s : P_C(s) \to P_D(f(s))$ acts on morphisms
        \emph{covariantly} (positions go \emph{forward});
  \item F0 says the codomain of $f^\flat_s(p)$ is $f(s \downarrow p)$;
  \item F1 says roots (identities) are preserved;
  \item F2 says shifts (compositions) are preserved, with a transport along F0.
\end{itemize}

The file also establishes that \lean{DContFunctor} forms a category (identity,
composition, and the three category laws, all by \lean{rfl}).

%===================================================================
\section{The \lean{DContFunctor} structure}
%===================================================================

% --- 1. DContFunctor ----------------------------------------------
\subsection{Structure: \lean{DContFunctor}}

\begin{defn}[Covariant functor of directed containers]
A \textbf{covariant directed-container functor} $F : C \to D$ consists of:
\begin{itemize}
  \item a \emph{covariant} shape map $f : S_C \to S_D$;
  \item a \emph{covariant} position map $f^\flat_s : P_C(s) \to P_D(f(s))$;
\end{itemize}
satisfying:
\begin{itemize}
  \item \textbf{F0:} $f(C \!\downarrow_s\! p) = D \!\downarrow_{f(s)}\! f^\flat_s(p)$
        (codomain preserved);
  \item \textbf{F1:} $f^\flat_s(o_C(s)) = o_D(f(s))$ (roots preserved);
  \item \textbf{F2:} $f^\flat_s(p \oplus_C q) = f^\flat_s(p) \oplus_D
        (\mathrm{F0} \triangleright f^\flat_{C \downarrow_s p}(q))$
        (shifts preserved, with transport along F0).
\end{itemize}
\end{defn}

\paragraph{Lean.}
\begin{lstlisting}[style=lean]
structure DContFunctor (C D : DirectedContainer) where
  onShapes : C.Shape -> D.Shape
  onPos : (s : C.Shape) -> C.Pos s -> D.Pos (onShapes s)
  f0 : forall (s : C.Shape) (p : C.Pos s),
    onShapes (C.sub s p) = D.sub (onShapes s) (onPos s p)
  f1 : forall s : C.Shape, onPos s (C.root s) = D.root (onShapes s)
  f2 : forall (s : C.Shape) (p : C.Pos s) (q : C.Pos (C.sub s p)),
    onPos s (C.shift s p q)
      = D.shift (onShapes s) (onPos s p)
          (f0 s p [> onPos (C.sub s p) q)
\end{lstlisting}

\paragraph{Haskell.}
\begin{lstlisting}[style=haskell]
data DContFunctor c d = DContFunctor
  { onShapes :: Shape c -> Shape d        -- covariant
  , onPos    :: Shape c -> Pos c -> Pos d -- covariant (positions go forward)
  -- F0: onShapes (sub_C s p) = sub_D (onShapes s) (onPos s p)
  -- F1: onPos s (root_C s) = root_D (onShapes s)
  -- F2: onPos s (shift_C s p q)
  --       = shift_D (onShapes s) (onPos s p)
  --           (onPos (sub_C s p) q)
  --   (transport along F0 is suppressed in Haskell)
  }
\end{lstlisting}

\begin{fidelitynote}
The key difference from \lean{DContMorphism}: \lean{onPos} maps $P_C(s)$ to
$P_D(f(s))$ (forward), whereas in a \lean{DContMorphism}, \lean{onPos} maps
$P_D(f(s))$ to $P_C(s)$ (backward).  In F2, the inner position
$f^\flat_{C \downarrow s}(q)$ lives over $f(C \!\downarrow_s\! p)$, but
$D.\mathrm{shift}$ requires it over $D \!\downarrow_{f(s)}\! f^\flat_s(p)$;
the transport $\mathrm{F0} \triangleright$ bridges these two equal shapes.
This transport is invisible in Haskell.
\end{fidelitynote}


%===================================================================
\section{Identity and composition}
%===================================================================

% --- 2. DContFunctor.id -------------------------------------------
\subsection{Definition: \lean{DContFunctor.id}}

\begin{defn}[Identity functor]
The identity covariant functor on $C$: both maps are $\mathrm{id}$.
All three laws hold by $\mathrm{rfl}$ (the F2 transport is along $\mathrm{rfl}$
and vanishes definitionally).
\end{defn}

\paragraph{Lean.}
\begin{lstlisting}[style=lean]
def DContFunctor.id (C : DirectedContainer) : DContFunctor C C where
  onShapes s := s
  onPos _ p := p
  f0 _ _ := rfl
  f1 _ := rfl
  f2 _ _ _ := rfl
\end{lstlisting}

\paragraph{Haskell.}
\begin{lstlisting}[style=haskell]
idDCF :: DContFunctor c c
idDCF = DContFunctor id (\_ p -> p)
\end{lstlisting}


% --- 3. onPos_transport -------------------------------------------
\subsection{Lemma: \lean{DContFunctor.onPos\_transport}}

\begin{prop}[Transport commutes with the position map]
Let $F : C \to D$ and let $e : a = b$ be a shape equality in $C$.  Then
\[
  F.\mathrm{onPos}_b(e \triangleright p)
  \;=\;
  \mathrm{congrArg}(F.\mathrm{onShapes}, e) \triangleright F.\mathrm{onPos}_a(p).
\]
That is, pushing a position forward after transporting along $e$ equals
transporting the image along the image equality.
\end{prop}

\paragraph{Lean.}
\begin{lstlisting}[style=lean]
theorem onPos_transport {C D : DirectedContainer} (phi : DContFunctor C D)
    {a b : C.Shape} (e : a = b) (p : C.Pos a) :
    phi.onPos b (e [> p) = congrArg phi.onShapes e [> phi.onPos a p
\end{lstlisting}

\paragraph{Haskell.}
\begin{lstlisting}[style=haskell]
-- Trivial in Haskell: transport is identity.
-- onPos b p = onPos a p  (since a = b implies same type)
\end{lstlisting}


% --- 4. transport_trans -------------------------------------------
\subsection{Lemma: \lean{DContFunctor.transport\_trans}}

\begin{prop}[Two transports equal one]
For any type family $P : \alpha \to \mathbf{Type}$, equalities $h_1 : x = y$
and $h_2 : y = z$, and element $p : P(x)$:
\[
  h_2 \triangleright (h_1 \triangleright p) \;=\; (h_1 \cdot h_2) \triangleright p.
\]
\end{prop}

\paragraph{Lean.}
\begin{lstlisting}[style=lean]
theorem transport_trans {alpha : Type} {P : alpha -> Type} {x y z : alpha}
    (h1 : x = y) (h2 : y = z) (p : P x) :
    h2 [> h1 [> p = (h1.trans h2) [> p
\end{lstlisting}

\paragraph{Haskell.}
\begin{lstlisting}[style=haskell]
-- Trivial in Haskell (transport is identity).
\end{lstlisting}


% --- 5. DContFunctor.comp -----------------------------------------
\subsection{Definition: \lean{DContFunctor.comp}}

\begin{defn}[Composition of covariant functors]
Given $F : C \to D$ and $G : D \to E$, their composite $G \circ F : C \to E$
has:
\begin{align*}
  \mathrm{onShapes}(s) &= G.\mathrm{onShapes}(F.\mathrm{onShapes}(s)), \\
  \mathrm{onPos}_s(p) &= G.\mathrm{onPos}_{F.\mathrm{onShapes}(s)}(F.\mathrm{onPos}_s(p)).
\end{align*}
Both maps compose forwards.  F0 is the composite of the factors' F0s; F1 uses
two rewrites; F2 uses both factors' F2s followed by the two transport lemmas.
\end{defn}

\paragraph{Lean.}
\begin{lstlisting}[style=lean]
def DContFunctor.comp {C D E : DirectedContainer}
    (phi : DContFunctor C D) (psi : DContFunctor D E) : DContFunctor C E where
  onShapes s := psi.onShapes (phi.onShapes s)
  onPos s p := psi.onPos (phi.onShapes s) (phi.onPos s p)
  f0 s p := (congrArg psi.onShapes (phi.f0 s p)).trans
    (psi.f0 (phi.onShapes s) (phi.onPos s p))
  f1 s := by rw [phi.f1, psi.f1]
  f2 s p q := by
    rw [phi.f2, psi.f2, psi.onPos_transport (phi.f0 s p), transport_trans]
\end{lstlisting}

\paragraph{Haskell.}
\begin{lstlisting}[style=haskell]
compDCF :: DContFunctor c d -> DContFunctor d e -> DContFunctor c e
compDCF phi psi = DContFunctor
  { onShapes = onShapes psi . onShapes phi
  , onPos    = \s p -> onPos psi (onShapes phi s) (onPos phi s p)
  }
\end{lstlisting}


%===================================================================
\section{Category laws for \lean{DContFunctor}}
%===================================================================

% --- 6. id_comp / comp_id / comp_assoc ----------------------------
\subsection{Theorems: identity and associativity}

\begin{thm}[Left unit]
For any $F : C \to D$, $\mathrm{comp}(\mathrm{id}_C, F) = F$ (by \lean{rfl}).
\end{thm}

\begin{thm}[Right unit]
For any $F : C \to D$, $\mathrm{comp}(F, \mathrm{id}_D) = F$ (by \lean{rfl}).
\end{thm}

\begin{thm}[Associativity]
For $F : B \to C$, $G : C \to D$, $H : D \to E$:
\[
  (H \circ G) \circ F \;=\; H \circ (G \circ F) \quad\text{(by \lean{rfl})}.
\]
\end{thm}

\paragraph{Lean.}
\begin{lstlisting}[style=lean]
theorem DContFunctor.id_comp {C D : DirectedContainer}
    (phi : DContFunctor C D) : comp (id C) phi = phi

theorem DContFunctor.comp_id {C D : DirectedContainer}
    (phi : DContFunctor C D) : comp phi (id D) = phi

theorem DContFunctor.comp_assoc {B C D E : DirectedContainer}
    (phi : DContFunctor B C) (psi : DContFunctor C D) (chi : DContFunctor D E) :
    comp (comp phi psi) chi = comp phi (comp psi chi)
\end{lstlisting}

\paragraph{Haskell.}
\begin{lstlisting}[style=haskell]
-- prop_id_comp phi = compDCF idDCF phi == phi
-- prop_comp_id phi = compDCF phi idDCF == phi
-- prop_assoc phi psi chi =
--   compDCF (compDCF phi psi) chi == compDCF phi (compDCF psi chi)
\end{lstlisting}


%===================================================================
\section{Translation conventions}
%===================================================================

\begin{enumerate}
\item \textbf{Covariant vs.\ contravariant.}  \lean{DContFunctor.onPos} maps
  $P_C(s) \to P_D(f(s))$ (forward); \lean{DContMorphism.onPos} maps
  $P_D(f(s)) \to P_C(s)$ (backward).  This is the defining distinction
  between functors and cofunctors of directed containers.

\item \textbf{F2 transport.}  The inner position in F2 must be transported
  along F0 before \lean{D.shift} can consume it.  All three category laws
  (\lean{rfl}) confirm that these transports compose cleanly.

\item \textbf{Design choice: direct definition.}  \lean{DContFunctor} is
  defined directly on containers (not via \lean{toSmallCat}), avoiding the
  round-trip gap and keeping transports $\mathrm{rfl}$-cheap for containers
  like \lean{network G} whose \lean{sub} is a projection.

\item \textbf{Proof bodies omitted.}  Only statements are translated.
\end{enumerate}

\end{document}
